% Teil II: durchgerechneter Weg, Teile a--d.

\subsection*{Rechnung}
\textbf{a) Periodendauer.} Die Reihe benutzt Basisfunktionen $\cos(nt)$,
$\sin(nt)$ mit ganzzahligem $n$; die Grundfrequenz ist $n=1$, also
\begin{equation}
  \boxed{\;T=2\pi\;}.
\end{equation}

\textbf{b) $f(t)=\cos(2t)\,\sin(3t)$.} Mit der Produkt-zu-Summe-Formel:
\begin{equation}
  \cos(2t)\sin(3t)=\tfrac12\big[\sin(3t+2t)-\sin(3t-2t)\big]
  =\tfrac12\sin(5t)-\tfrac12\sin(t)\;.
\end{equation}
Jeder Term ist bereits eine Basisfunktion:
\begin{equation}
  a_0=a_n=0,\qquad
  \boxed{\;b_1=-\tfrac12,\quad b_5=+\tfrac12\;},\quad
  \text{alle übrigen }b_n=0\;.
\end{equation}

\textbf{c) $f(t)=\cos(2t)+\sin(3t)$.} Direkt ablesbar:
\begin{equation}
  \boxed{\;a_2=1,\qquad b_3=1\;},\qquad
  a_0=0,\; \text{alle anderen }a_n=b_n=0\;.
\end{equation}

\textbf{d) $f(t)=\cos(2t)+\cos(t/\pi)$.}
$\cos(2t)$ liefert sofort $a_2=1$. Der Term $\cos(t/\pi)$ hat die Periode
$2\pi^2$: Das Verhältnis zur Grundperiode $2\pi$ ist $\pi$ (irrational),
die Frequenz $1/\pi$ ist \emph{kein ganzzahliges Vielfaches} von $1$. Deshalb
lässt sich $\cos(t/\pi)$ durch die $2\pi$-Basis \emph{nicht endlich} darstellen;
seine Anteile streuen auf alle Koeffizienten:
\begin{equation}
  a_n^{\,\mathrm{inkom}}
  =\frac1\pi\int_0^{2\pi}\cos\!\big(\tfrac{t}{\pi}\big)\cos(nt)\,\mathrm dt
  \neq 0 \quad(\forall n),\qquad
  b_n=0\;.
\end{equation}
Zusammengefasst:
\begin{equation}
  \boxed{\;a_2=1,\quad a_n=\frac1\pi\int_0^{2\pi}\cos\!\big(\tfrac{t}{\pi}\big)
  \cos(nt)\,\mathrm dt\ \text{für alle }n\neq2,\quad a_0=0,\ b_n=0\;}.
\end{equation}
Die eigentliche Aussage von (d): Nicht jede Funktion ist durch eine
gegebene Fourier-Basis mit \emph{ganzzahligen} Frequenzen darstellbar ---
frequenzinkommensurable Anteile brauchen unendlich viele (kleine) Terme.

\begin{ergebnis}
(a) $T=2\pi$; (b) $b_1=-\tfrac12$, $b_5=+\tfrac12$;
(c) $a_2=1$, $b_3=1$; (d) $a_2=1$, der $\cos(t/\pi)$-Anteil ist inkommensurabel
und für alle $n$ über die angegebenen Integrale bestimmt.
\end{ergebnis}

\begin{kontrolle}
\textbf{Querprobe (b):} Rückeinsetzen $-\tfrac12\sin t+\tfrac12\sin5t
=\tfrac12(\sin5t-\sin t)=\cos2t\,\sin3t$ --- korrekt. \checkmark

\textbf{(c):} Die Summe ist bereits ihre eigene Reihe --- konsistent mit
$a_2=1$, $b_3=1$. \checkmark

\textbf{(d):} Mit $n=1$ ist $\cos t\cos(t/\pi)$ über $[0,2\pi]$ nicht
orthogonal zu null, also $a_1\neq0$ --- wie erwartet für eine inkommensurable
Frequenz.
\end{kontrolle}

\begin{merksatz}
Wenn $f$ bereits als Summe von Basisfunktionen $\cos(nt)$/$\sin(nt)$ vorliegt,
liest man die Koeffizienten direkt ab; nur Frequenzen, die zur Periode
inkommensurabel sind (kein ganzzahliges Vielfaches der Grundfrequenz),
machen echte Integral-Rechnung nötig.
\end{merksatz}
